% lower_bound.tex - Lower bound on regret
\documentclass[11pt]{article}
\usepackage{amsmath, amssymb, amsthm}

\title{A Lower Bound for Adversarial Spread Selection}
\author{}
\date{}

\newtheorem{theorem}{Theorem}

\begin{document}
\maketitle

\section{Statement}
\begin{theorem}
For $K\ge 2$ arms and horizon $T$, every bandit algorithm has an adversarial
reward sequence with expected regret
\[
\mathbb{E}[R_T] = \Omega(\sqrt{TK}).
\]
\end{theorem}

\noindent\textbf{Note.} The $\Omega(\sqrt{TK})$ lower bound does not include the
$\sqrt{\log K}$ factor that appears in Exp3's upper bound. A tighter
construction for the oblivious adversary model yields
$\Omega(\sqrt{TK\log K})$ (see Auer et al., 2002, Theorem 3.3).

\section{Proof Sketch}
It is enough to consider stochastic reward sequences chosen before play. Pick
one arm $a$ uniformly from $[K]$. Arm $a$ has Bernoulli mean $1/2+\epsilon$ and
all other arms have mean $1/2$. If the learner pulls arm $a$ exactly $N_a$
times, then its regret is
\[
\epsilon\,\mathbb{E}[T-N_a].
\]
Identifying $a$ requires distinguishing $K$ nearby Bernoulli distributions.
Standard change-of-measure arguments imply that, for
$\epsilon = c\sqrt{K/T}$ with a small universal constant $c$, the learner cannot
place more than a constant fraction of its pulls on the hidden best arm in
expectation. Hence
\[
\mathbb{E}[T-N_a] = \Omega(T),
\qquad
\mathbb{E}[R_T] = \Omega(T\epsilon)=\Omega(\sqrt{TK}).
\]
Yao's minimax principle transfers the distributional construction to an
adversarial lower bound.

\end{document}
